\documentclass[11pt,a4paper]{article}
\usepackage[T1]{fontenc}
\usepackage[utf8]{inputenc}
\usepackage{lmodern}
\usepackage[margin=28mm]{geometry}
\usepackage{amsmath,amssymb,amsthm}
\usepackage{microtype,booktabs,cite}
\usepackage[hidelinks,breaklinks]{hyperref}
\hypersetup{
  pdftitle={Exact directional NFA separation for binary words with one marker},
  pdfauthor={Xiaojian Zeng},
  pdfkeywords={word separation, nondeterministic finite automata, Boolean relations}
}
\newtheorem{theorem}{Theorem}
\newtheorem{lemma}[theorem]{Lemma}
\newtheorem{proposition}[theorem]{Proposition}
\newtheorem{corollary}[theorem]{Corollary}
\DeclareMathOperator{\nsep}{nsep}
\DeclareMathOperator{\lcm}{lcm}
\newcommand{\Aone}{A_{\mathrm{one}}}
\newcommand{\Nzero}{\mathbb Z_{\geq 0}}
\title{Exact directional NFA separation\\for binary words with one marker}
\author{Xiaojian Zeng\\
  \small PBC School of Finance, Tsinghua University, 100190\\
  \small \href{mailto:xiaojianzeng@foxmail.com}{xiaojianzeng@foxmail.com}}
\date{}
\begin{document}
\maketitle

\begin{abstract}
For distinct words \(w,x\), let \(\nsep(w,x)\) be the least number of
states in a nondeterministic finite automaton accepting \(w\) and rejecting
\(x\). We determine the largest possible value of this directional cost
over equal-length binary words containing exactly one \(1\), when the
opposite cost is at most \(q\). For every integer \(q\geq3\), the maximum is
\((q-1)^2+2\), and it is attained. The upper bound combines the classical
index bound for powers of binary relations with explicit separators.
Sharpness follows by extending Nicol's cycle-insertion argument to an
initial path that visits every state. Thus the entire one-marker class
has a sharp quadratic bound. The asymptotic growth of the unrestricted
directional extremal function remains unresolved.
\end{abstract}

\section{The restricted extremal problem}

The word-separation problem asks how many states are needed to distinguish
two specified words; its deterministic formulation appears in
Goralčík and Koubek~\cite{goralcik1986}.
Demaine, Eisenstat, Shallit and Wilson~\cite[Section~6]{demaine2011}
introduced the directional quantity \(\nsep(w,x)\), exhibited its
asymmetry, and asked whether the difference between the two orientations
could be unbounded. Nicol~\cite[Theorem~6]{nicol2026} answered this question
by constructing equal-length words with
\(\nsep(u_s,v_s)>s^2\) and \(\nsep(v_s,u_s)\leq s+1\), for every \(s\geq2\).

Throughout this note an NFA has a finite nonempty state set, alphabet
\(\{0,1\}\), one initial state, any subset of accepting states, and a
possibly partial transition relation. There are no \(\varepsilon\)-edges.
Acceptance means the existence of a path from the initial state to an
accepting state with the specified label. Only the acceptance of \(w\)
and rejection of \(x\) enter the definition of \(\nsep(w,x)\); behavior on
every other word is unrestricted. These are the conventions of
\cite[Section~4]{nicol2026}.

Nicol~\cite[Section~4.1]{nicol2026} defines
\[
 A(q)=\max\{\nsep(w,x): w,x\in\{0,1\}^*,\
       |w|=|x|,\ w\ne x,\ \nsep(x,w)\leq q\}
\]
and records the bounds
\begin{equation}\label{eq:prior}
 (q-1)^2+1\leq A(q)\leq 2^q,\qquad q\geq3.
\end{equation}
The upper bound follows by determinizing a reverse separator by the
subset construction~\cite[Definition~11 and Theorem~11]{rabin1959}
and complementing it. The deterministic states are all subsets, including
the empty subset, so the machine is complete and uses at most \(2^q\)
states. This bound also ensures that the indicated integer maximum exists.
The asymptotic growth question for \(A(q)\) remains unresolved here.

We restrict both words to the language \(0^*10^*\). Explicitly, put
\[
 \Aone(q)=\max\left\{
 \begin{array}{@{}l@{}}
 \nsep(0^a1\,0^b,0^c1\,0^d):\
 a,b,c,d\in\Nzero,\ a+b=c+d,\ a\ne c,\\[2pt]
 \hfill \nsep(0^c1\,0^d,0^a1\,0^b)\leq q
 \end{array}\right\}.
\]
All lengths are allowed, and the four zero blocks may be empty.

\begin{theorem}\label{thm:main}
For every integer \(q\geq3\),
\[
 \Aone(q)=(q-1)^2+2.
\]
The maximum is attained by a finite pair of equal-length words.
\end{theorem}

The restriction contains Nicol's quadratic witness family. The main
contribution is a matching upper bound for every admissible one-marker
pair, together with a sharpening of that family's hard-direction cost
by one state. The consecutive-cycle constructions and unary-block
indistinguishability arguments in
\cite[Theorems~3--4]{demaine2011} are antecedents; those results do not
bound one directional NFA cost as a function of the other.
Neither \cite{demaine2011} nor \cite{nicol2026} states the restricted
maximum in Theorem~\ref{thm:main}. The comparison is with these inspected
sources, rather than a claim of exhaustive bibliographic priority.

\begin{corollary}\label{cor:general}
For every \(q\geq3\), \(A(q)\geq(q-1)^2+2\).
No pair of equal-length words in \(0^*10^*\) can require more than
\((q-1)^2+2\) states in one orientation while requiring at most \(q\)
states in the other.
\end{corollary}

\section{Powers of the zero relation}

For a binary relation \(R\) on \(n\) points, \(R^t\) records the endpoints
of walks of length exactly \(t\); \(R^0\) is the identity relation.
Composition is equivalently Boolean matrix multiplication, using OR and
AND. Write
\[
 H_n=(n-1)^2+1,\qquad L_n=\lcm(1,\ldots,n).
\]
We use the following classical index theorem.

\begin{proposition}[Schwarz's index bound]\label{prop:schwarz}
If \(n>1\), the powers of every binary relation on \(n\) points are
periodic from some index at most \(H_n\).
\end{proposition}

The precise all-relations statement is given in
Schwarz~\cite[p.~167, proof of Theorem~3, Case~1]{schwarz1974}:
the power \(R^{n^2-2n+2}\) lies in the periodic group of powers, defined
on p.~161. This is the author's published account of an earlier result;
we invoke the statement in that account. It applies to reducible as well
as irreducible relations. The same index bound is reported in
\cite[Theorem~1]{zajac2012}. We give the period consequence needed below
directly, so no separate formula for the universal period is required.

\begin{lemma}\label{lem:period}
For \(n>1\), every binary relation \(R\) on \(n\) points satisfies
\begin{equation}\label{eq:period}
 R^{t+L_n}=R^t\qquad(t\geq H_n).
\end{equation}
\end{lemma}
\begin{proof}
A walk of length \(t\geq n\) has a repeated vertex and contains a simple
directed cycle of length \(h\leq n\). Inserting \(L_n/h\) additional
traversals of this cycle preserves the endpoints and increases the
length by \(L_n\). Thus \(R^t\subseteq R^{t+L_n}\) whenever \(t\geq n\).

Since \(H_n\geq n\), fix \(t\geq H_n\), and let \(e\geq1\) be a period
valid from \(H_n\) onward, supplied by Proposition~\ref{prop:schwarz}.
Then
\[
 R^t\subseteq R^{t+L_n}\subseteq\cdots
       \subseteq R^{t+eL_n}=R^t.
\]
Every inclusion is therefore equality.
\end{proof}

Consequently, two exponents at least \(H_n\) that are congruent modulo
\(L_n\) give equal powers. An NFA with fewer than \(n\) states can be
padded by unreachable isolated states when applying this identity.

\section{The uniform upper bound}

Fix \(q\geq3\) and an admissible ordered pair
\[
 w=0^a1\,0^b,\qquad x=0^c1\,0^d.
\]
Set \(H=H_q\), \(L=L_q\), and \(\Delta=a-c\ne0\).
Equal length gives \(b-d=-\Delta\). By admissibility some NFA with at
most \(q\) states accepts \(x\) and rejects \(w\).

\paragraph{A separating modulus.}
Suppose \(L\nmid\Delta\). Then some \(r\in\{2,\ldots,q\}\) does not
divide \(\Delta\). Take a directed zero-cycle on residues
\(0,\ldots,r-1\), starting at \(0\). Add a separate accepting state \(f\)
with a zero-loop, and one \(1\)-edge, from residue \(a\bmod r\) to \(f\).
There are no other edges, and \(f\) is the sole accepting state.
The machine accepts \(w\). It rejects \(x\), since the marker in \(x\)
is read at a different residue. The state count is
\[
 r+1\leq q+1\leq H+1.
\]

\paragraph{A short block.}
Suppose one of \(a,b,c,d\) is less than \(H\). First let
\(k=\min(a,c)<H\). The words first differ immediately after the common
prefix \(0^k\).

If the word to be accepted has the earlier marker, use a zero-chain
of \(k\) edges from the initial state to a state \(p\), a \(1\)-edge
from \(p\) to a new accepting state \(f\), and a zero-loop at \(f\).
All remaining transitions are absent. The later-marker word fails
on its next zero at \(p\).
If the word to be accepted has the later marker, use a zero-chain
of \(k+1\) edges and put both symbol loops at its final state, which
alone is accepting. The earlier-marker word fails on its \(1\).
Both constructions have \(k+2\leq H+1\) states.

If instead \(k=\min(b,d)<H\), apply these constructions to the two
reversed words. Each constructed machine has a single accepting state.
Reverse every edge and exchange its initial and accepting states.
Path reversal proves that the resulting machine accepts \(w\) and
rejects \(x\), with the same number of states and one initial state.
This is the path-reversal principle used in
\cite[Theorem~5]{demaine2011}, applied here to the explicit machines.
Empty prefixes and suffixes are included.

\paragraph{The remaining case is impossible.}
It remains to consider \(L\mid\Delta\) and \(a,b,c,d\geq H\).
Pad the assumed reverse separator to \(q\) states, and let \(B,C\)
be its zero and one relations. Lemma~\ref{lem:period} gives
\[
 B^a=B^c,\qquad B^b=B^d,\qquad
 B^aCB^b=B^cCB^d.
\]
The two words induce the same relation on the state set. Hence no
initial state and accepting subset can distinguish them, contradicting
the reverse separator.

These cases exhaust the possibilities. Every admissible pair therefore
has \(\nsep(w,x)\leq H+1\), proving
\begin{equation}\label{eq:upper}
 \Aone(q)\leq(q-1)^2+2.
\end{equation}

\section{An attaining pair}

Fix \(q\geq3\) and set
\begin{equation}\label{eq:parameters}
 s=q-1,\qquad M=s^2,\qquad T=2^{M+1},
 \qquad P=\lcm(1,\ldots,T).
\end{equation}
Consider
\begin{equation}\label{eq:witness}
 u=0^M1\,0^{T+P},\qquad v=0^{M+P}1\,0^T.
\end{equation}
The words are distinct and have common length \(M+T+P+1\).
This adapts the pair in \cite[Theorem~6]{nicol2026}, replacing
its value \(T=2^M\) by \(2^{M+1}\). The additional step below handles
a prefix that has not repeated a vertex but has visited every state.

\begin{lemma}\label{lem:hard}
For the pair in \eqref{eq:witness}, \(\nsep(u,v)=M+2\).
\end{lemma}
\begin{proof}
Suppose an \(m\)-state NFA, with \(m\leq M+1\), accepts \(u\).
Choose an accepting path and mark its states adjacent to the marker:
\[
 i\overset{0^M}{\longrightarrow}p\overset{1}{\longrightarrow}z
      \overset{0^{T+P}}{\longrightarrow}f.
\]
If the first segment repeats a vertex, it contains a simple zero-cycle.
Otherwise it visits \(M+1\) different vertices, forcing \(m=M+1\)
and visiting every state. The final segment has more than \(m\) edges,
so it contains a simple zero-cycle, one of whose vertices also occurs
on the first segment. In either case some vertex of the prefix supports
a zero-cycle of length \(h\leq m\).
Since \(h\leq M+1\leq T\), we have \(h\mid P\).
Inserting \(P/h\) traversals of that cycle at its prefix vertex
gives a path \(i\overset{0^{M+P}}{\longrightarrow}p\).
The cycle need not have been traversed in the original prefix.

Let \(S_j\) be the subset reachable from \(z\) after exactly \(j\) zeros.
The update of these subsets is deterministic and there are at most
\(2^m\leq T\) possible subsets. Thus \(S_\alpha=S_\beta\) for some
\(0\leq\alpha<\beta\leq T\). The sequence is periodic from \(\alpha\)
onward with period \(c=\beta-\alpha\leq T\).
As \(c\mid P\) and \(T\geq\alpha\), this implies
\[
 S_T=S_{T+P}.
\]
In particular, the chosen final state \(f\) is reachable from \(z\)
after \(T\) zeros. Joining this shorter suffix to the extended prefix
and the same \(1\)-edge produces an accepting path on \(v\).
No NFA with at most \(M+1\) states can therefore accept \(u\) and reject \(v\).

Conversely, the earlier-marker chain from the upper-bound construction,
with \(k=M\), accepts \(u\) and rejects \(v\) using \(M+2\) states.
\end{proof}

\begin{lemma}\label{lem:easy}
For the same pair, \(\nsep(v,u)\leq q\).
\end{lemma}
\begin{proof}
Use Nicol's graph~\cite[Figure~1]{nicol2026} on states \(0,1,\ldots,s\).
Its zero-edges are
\[
 j\longrightarrow j+1\quad(0\leq j<s),\qquad
 s\longrightarrow0,\quad s\longrightarrow1.
\]
The only \(1\)-edge is a loop at \(0\), which is both initial and
the sole accepting state.

The positive lengths of zero-walks from \(0\) back to \(0\) are exactly
\begin{equation}\label{eq:returns}
 \{\alpha s+\beta(s+1):\alpha,\beta\in\mathbb Z,\
      \alpha\geq0,\ \beta\geq1\}.
\end{equation}
Indeed, each return traverses the \((s+1)\)-cycle, with any number of
traversals of the \(s\)-cycle through \(1,\ldots,s\) inserted before
returning; concatenating returns gives \eqref{eq:returns}.
Every indicated combination can be realized.

Represent residues modulo \(s\) by \(r=1,\ldots,s\).
The least positive return length in residue \(r\) is \(r(s+1)\);
all later lengths in that residue are obtained by adding multiples
of \(s\). Therefore the largest excluded positive length is
\[
 \max_{1\leq r\leq s}\{r(s+1)-s\}=s^2=M.
\]
The empty return is also possible and does not affect this calculation.
Thus \(M\) is not a return length, while every length greater than \(M\)
is. The machine rejects \(u\), since after its first \(M\) zeros no path
has reached the sole available marker transition.
It accepts \(v\), since \(M+P>M\) and \(T>M\) are both return lengths.
The graph has \(s+1=q\) states.
\end{proof}

Lemmas~\ref{lem:hard} and \ref{lem:easy} supply an admissible pair
attaining \((q-1)^2+2\). Together with \eqref{eq:upper} they prove
Theorem~\ref{thm:main} and Corollary~\ref{cor:general}.
The attainment argument is elementary and does not use
Proposition~\ref{prop:schwarz}.

\section{Exact diagnostics and limitations}\label{sec:checks}

The supplement checks relation powers by finite orbit indexing and
simulates the explicit machines by exact subset propagation.
Enumerating Boolean-matrix indices and periods is a classical
computational problem studied by Zajac and J\'okay~\cite{zajac2012};
our enumeration is smaller and also checks the mixed-symbol transfer
used in the lower-bound argument.

Two features matter for these checks: exact local predicates and explicit
coverage of the enumerated domain. A related verification design appears
in Zeng et al.~\cite[Sections~3, 4 and 6]{zeng2026finite}, where finite
graph certificates are checked using integer predicates, reconstructed
orbit coverage, and a separate list-based implementation alongside a
bitset implementation. Here the development checks use Boolean row
bitsets, while the separately written supplementary verifier uses sets
of destinations and exact orbit tables. The connection is methodological;
no geometric theorem from that paper is used in our proof.

The supplementary run completed with no failed assertions.
Its exhaustive relation domains are shown in Table~\ref{tab:coverage}.
For every unary relation \(B\) it checks the index/period identity and
\(B^T=B^{T+P}\), where \(T=2^n\) and \(P=\lcm(1,\ldots,T)\).
For \(n\leq3\), it additionally enumerates every second relation \(C\)
and checks the entrywise inclusion
\begin{equation}\label{eq:diagnostic}
 B^{n-1}CB^{T+P}\subseteq B^{n-1+P}CB^T.
\end{equation}
Entrywise inclusion covers every initial state and accepting subset.
There is no complete binary-transition census for \(n=4\).

\begin{table}[ht]
\centering
\caption{Exhaustive domains in the supplementary verifier.}\label{tab:coverage}
\begin{tabular}{@{}rrr@{}}
\toprule
States \(n\) & Unary relations \(B\) & Binary pairs \((B,C)\)\\
\midrule
1 & 2 & 4\\
2 & 16 & 256\\
3 & 512 & 262,144\\
4 & 65,536 & not enumerated\\
\bottomrule
\end{tabular}
\end{table}

The construction tests cover \(q=3,\ldots,12\): 16,480 prefix machines,
16,480 reversed machines, and 8,000 modular machines.
They include all short-block thresholds, both orientations, and empty
blocks in the specified test families. Exact compressed witnesses for
\(q=3,4\) are also checked. For example, \(q=3\) gives
\(M=4\), \(T=32\), and \(P=144403552893600\).
The checks verify the two explicit machines' behavior, not the six-state
minimum; Lemma~\ref{lem:hard} supplies that minimum.
As a negative control, a chain disjoint from a loop shows the failure
of the overstrong insertion assertion at prefix length \(n-2\),
for each \(n=2,\ldots,12\).
The supplementary instructions specify the full test families.

These finite computations are diagnostics for the symbolic proof.
They do not prove the statement for untested numbers of states.
The exact restricted optimum leaves the gap between quadratic and
exponential growth in the unrestricted problem. In particular,
any superquadratic lower-bound construction for \(A(q)\) must use
pairs outside \(0^*10^*\). We make no claim about the shortest attaining
word length, the least reverse cost of the displayed witnesses,
or a classification of all attaining pairs.

\paragraph{Computational and AI assistance.}
OpenAI Codex assisted with literature searches, proof development,
exact verification code, a separate mathematical review pass, and
manuscript preparation. The separate review pass was performed in a
fresh model context and was not external peer review.
The accompanying Python~3.10-compatible verifier uses only the standard
library and exact finite-set and integer operations; no floating-point
tolerance or random sampling is used.
No proof assistant was used. The proof above and its stated classical
dependency, rather than model output or a successful computation,
support the all-size theorem.

\bibliographystyle{plainurl}
\bibliography{references}
\end{document}
